
%(BEGIN_QUESTION)
% Copyright 2006, Tony R. Kuphaldt, released under the Creative Commons Attribution License (v 1.0)
% This means you may do almost anything with this work of mine, so long as you give me proper credit

I 1808 publiserte den engelske læreren John Dalton flere postulater om stoffer, som representerte datidens atomteori:

\begin{itemize}
\item{} Et grunnstoff består av svært små udelelige partikler kalt atomer.
\item{} Alle atomer av et gitt grunnstoff har identiske egenskaper, som skiller seg fra andre grunnstoffers egenskaper.
\item{} Atomer kan ikke skapes, ødelegges eller omdannes til atomer av et annet grunnstoff.
\item{} Forbindelser dannes når atomer av ulike grunnstoffer kombineres i små heltallsforhold.
\item{} Relativt antall og type atomer er konstant i en gitt forbindelse.
\end{itemize} 

Avgjør om disse postulatene fortsatt stemmer med moderne atomteori, og beskriv eventuelle avvik du finner.

\underbar{file i00556}
%(END_QUESTION)





%(BEGIN_ANSWER)

{\bf Et grunnstoff består av svært små udelelige partikler kalt atomer.} Dette er i tråd med moderne teori, med unntak av at vi nå vet at atomer {\it kan} deles i mindre komponenter (elementærpartikler), som igjen kan deles videre (kvarker osv.).

\vskip 10pt

{\bf Alle atomer av et gitt grunnstoff har identiske egenskaper, som skiller seg fra andre grunnstoffers egenskaper.} Isotoper har {\it litt} ulike egenskaper innenfor samme grunnstoff. For noen grunnstoffer som hydrogen (H) er isotop-masseforholdene betydelige (Tritium:Deuterium:Hydrogen = 3:2:1). For de fleste isotoper er masseforskjellen liten og krever presisjonsutstyr for å måles. På Daltons tid fantes ikke slikt utstyr.

\vskip 10pt

{\bf Atomer kan ikke skapes, ødelegges eller omdannes til atomer av et annet grunnstoff.} Moderne vitenskap sier at {\it materie} ikke kan skapes eller ødelegges, men kan omformes mellom ulike former (inkludert energi). Atomer kan definitivt {\it transmuteres} til andre grunnstoffer. I naturen skjer dette i liten skala, og på større skala kreves avansert utstyr som ikke var tilgjengelig på Daltons tid.

\vskip 10pt

{\bf Forbindelser dannes når atomer av ulike grunnstoffer kombineres i små heltallsforhold.} Dette er helt i samsvar med moderne atomteori, og er en god definisjon av en ``forbindelse''.

\vskip 10pt

{\bf Relativt antall og type atomer er konstant i en gitt forbindelse.} Også dette er fullt i samsvar med moderne teori.

%(END_ANSWER)





%(BEGIN_NOTES)


%INDEX% Kjemi, grunnprinsipper: Daltons lover

%(END_NOTES)

